% chapters/04/section-01.tex
\section{Kerangka Siklus Permainan}
\label{sec:kerangka-siklus}

Bab~II telah menjelaskan mengapa program permainan berjalan terus dan
mengapa waktu harus dikelola. Bab ini mewujudkannya menjadi kode. Bagian
pertama membahas kerangkanya; bagian-bagian berikutnya membahas setiap tahap
secara terperinci.

\subsection{Tiga Tahap yang Selalu Sama}

Siklus permainan terdiri atas tiga tahap yang dikerjakan berurutan dan
diulang terus-menerus, seperti diperlihatkan pada
Gambar~\ref{fig:siklus-game-loop}.

\begin{figure}[htbp]
\centering
% chapters/04/figures/siklus-game-loop.tikz
% Diagram siklus input -> update -> render untuk Bab IV.
% Dipanggil dari section-*.tex dengan % chapters/04/figures/siklus-game-loop.tikz
% Diagram siklus input -> update -> render untuk Bab IV.
% Dipanggil dari section-*.tex dengan % chapters/04/figures/siklus-game-loop.tikz
% Diagram siklus input -> update -> render untuk Bab IV.
% Dipanggil dari section-*.tex dengan \input{figures/siklus-game-loop.tikz}
\begin{tikzpicture}[
  node distance=20mm,
  tahap/.style={
    rectangle, rounded corners=2pt, draw=ubPrimary, thick, fill=ubAlt,
    text width=27mm, align=center, minimum height=12mm
  },
  panah/.style={-{Stealth[length=2.6mm]}, thick, draw=ubAccent}
]
\node[tahap] (input) {\small\bfseries Input\\[1pt]\footnotesize baca peristiwa};
\node[tahap, right=of input] (update) {\small\bfseries Update\\[1pt]\footnotesize perbarui keadaan};
\node[tahap, right=of update] (render) {\small\bfseries Render\\[1pt]\footnotesize gambar layar};

\draw[panah] (input) -- (update);
\draw[panah] (update) -- (render);
\draw[panah] (render.south) -- ++(0,-7mm) -| (input.south);
\node[below=11mm of update, font=\footnotesize\itshape, text=ubPrimary]
  {diulang terus-menerus selama permainan berjalan};
\end{tikzpicture}

\begin{tikzpicture}[
  node distance=20mm,
  tahap/.style={
    rectangle, rounded corners=2pt, draw=ubPrimary, thick, fill=ubAlt,
    text width=27mm, align=center, minimum height=12mm
  },
  panah/.style={-{Stealth[length=2.6mm]}, thick, draw=ubAccent}
]
\node[tahap] (input) {\small\bfseries Input\\[1pt]\footnotesize baca peristiwa};
\node[tahap, right=of input] (update) {\small\bfseries Update\\[1pt]\footnotesize perbarui keadaan};
\node[tahap, right=of update] (render) {\small\bfseries Render\\[1pt]\footnotesize gambar layar};

\draw[panah] (input) -- (update);
\draw[panah] (update) -- (render);
\draw[panah] (render.south) -- ++(0,-7mm) -| (input.south);
\node[below=11mm of update, font=\footnotesize\itshape, text=ubPrimary]
  {diulang terus-menerus selama permainan berjalan};
\end{tikzpicture}

\begin{tikzpicture}[
  node distance=20mm,
  tahap/.style={
    rectangle, rounded corners=2pt, draw=ubPrimary, thick, fill=ubAlt,
    text width=27mm, align=center, minimum height=12mm
  },
  panah/.style={-{Stealth[length=2.6mm]}, thick, draw=ubAccent}
]
\node[tahap] (input) {\small\bfseries Input\\[1pt]\footnotesize baca peristiwa};
\node[tahap, right=of input] (update) {\small\bfseries Update\\[1pt]\footnotesize perbarui keadaan};
\node[tahap, right=of update] (render) {\small\bfseries Render\\[1pt]\footnotesize gambar layar};

\draw[panah] (input) -- (update);
\draw[panah] (update) -- (render);
\draw[panah] (render.south) -- ++(0,-7mm) -| (input.south);
\node[below=11mm of update, font=\footnotesize\itshape, text=ubPrimary]
  {diulang terus-menerus selama permainan berjalan};
\end{tikzpicture}

\caption{Siklus permainan: tiga tahap yang diulang setiap bingkai}
\label{fig:siklus-game-loop}
\end{figure}

\begin{enumerate}
  \item \textbf{Masukan} (\istilah{input}). Membaca peristiwa yang terjadi sejak
        bingkai sebelumnya: tombol yang ditekan, tetikus yang digeser, atau
        jendela yang hendak ditutup. Hasil tahap ini adalah gambaran tentang
        apa yang sedang diinginkan pemain.
  \item \textbf{Pembaruan} (\istilah{update}). Mengubah keadaan permainan
        berdasarkan masukan tersebut dan berdasarkan berjalannya waktu.
        Pekerjaan pada tahap ini meliputi pemindahan objek, pemeriksaan
        tabrakan, penambahan skor, dan penentuan perilaku objek yang
        dikendalikan komputer.
  \item \textbf{Penayangan} (\istilah{render}). Menggambar ulang seluruh isi
        jendela menurut keadaan yang baru saja diperbarui, lalu
        menampilkannya kepada pemain.
\end{enumerate}

Ketiga tahap itu berlangsung dalam waktu yang sangat singkat, biasanya jauh
di bawah sepersepuluh detik, lalu diulang. Karena pengulangannya cepat,
pemain melihat satu gambar yang terus berubah dan merasakannya sebagai gerak
yang hidup.

\subsection{Mengapa Urutannya Demikian}

Urutan tahapan bukan pilihan gaya penulisan, melainkan keharusan yang lahir
dari sifat pekerjaannya.

\begin{itemize}
  \item \textbf{Masukan didahulukan} karena pembaruan memerlukan pengetahuan
        tentang kehendak pemain. Apabila pembaruan dijalankan lebih dahulu,
        tombol yang baru ditekan baru akan berpengaruh pada bingkai berikutnya,
        dan pemain merasakan tanggapan yang terlambat satu bingkai.
  \item \textbf{Penayangan dikerjakan terakhir} karena yang seharusnya terlihat
        pemain adalah keadaan akhir setelah seluruh perhitungan selesai.
        Menggambar di tengah perhitungan berarti menampilkan keadaan setengah
        jadi.
\end{itemize}

Kedua alasan itu kedengaran sepele pada permainan sederhana, dan memang
demikian. Namun ketika jumlah objek bertambah dan sebagian perhitungan saling
bergantung, urutan yang tidak konsisten menimbulkan gejala yang sulit
dilacak: objek yang tampak bergerak mendahului tabrakannya, atau skor yang
bertambah sebelum kejadiannya sungguh terjadi.

\begin{catatabox}
Pada Bab~XIII, ketika Anda beralih ke Unity, ketiga tahap itu tetap ada
walaupun sebagian dikelola mesin. Unity menyediakan tahap tersendiri untuk
menjalankan perintah Anda setiap bingkai, dan tahap itu setara dengan tahap
pembaruan di sini. Yang berubah hanyalah siapa yang menulis perulangannya.
\end{catatabox}

\subsection{Kerangka dalam Bentuk Kode}

Kerangka siklus permainan, tanpa isi setiap tahapnya, berbentuk sesederhana
berikut.

\begin{kodekeluaran}
running = True

while running:
    # Tahap 1: baca peristiwa masukan
    # Tahap 2: perbarui keadaan permainan
    # Tahap 3: gambar ulang seluruh jendela
\end{kodekeluaran}

Perhatikan tiga hal pada kerangka tersebut.

\begin{enumerate}
  \item \textbf{Perulangannya dikendalikan sebuah variabel}, pada contoh ini
        \code{running}. Variabel itu diubah menjadi salah ketika pemain
        meminta keluar, sehingga siklus berakhir dengan rapi. Perulangan tanpa
        cara keluar hanya dapat dihentikan dengan mematikan program secara
        paksa.
  \item \textbf{Variabel keadaan berada di luar perulangan.} Keadaan permainan
        harus tetap hidup dari satu bingkai ke bingkai berikutnya. Apabila
        variabel keadaan dibuat di dalam perulangan, isinya kembali ke nilai
        awal setiap bingkai dan permainan kehilangan ingatannya.
  \item \textbf{Ketiga tahap berada di dalam satu perulangan yang sama.} Ketiga
        tahap itu bukan tiga perulangan terpisah; pemisahan menjadi perulangan
        tersendiri membuat tahap-tahap itu berjalan tidak serempak, dan
        perilaku permainan menjadi sulit diperkirakan.
\end{enumerate}

Kerangka inilah yang akan Anda lengkapi pada bagian-bagian berikutnya,
sehingga menjadi siklus permainan yang benar-benar berjalan
\cite{nystrom-gameloop,sweigart-mg}.
